These three sonic sculptures—\emph{Buffalo Soldier}, \emph{Did Sun Ra Fear Cryosleep?}, and \emph{Don't Play with My Hair}—collectively demonstrate a practice of culturally grounded technological modification that generates knowledge through making. Each work investigates different material substrates (earth and rock, water and sediment, hair and gesture) and employs different technical approaches (piezoelectric sensing and spatialization, hydrophone sonification and vertical audio, computer vision and synthesis), yet all three share a fundamental commitment: they refuse the separation between cultural speculation and technological practice, insisting that Black epistemologies can and must guide how tools are modified, how parameters are mapped, and how sonic experiences are structured. It also does feel as if, through the iterative development of these works, a recurring set of design principles has begun to surface: a synthesis methodology grounded more in the direct translation of material affordances and properties, gestural vocabularies, and ritualized interactions. This approach seems, in many ways, to stand in productive tension with more well-known synthesis models that prioritize mathematical abstraction.

As research instruments, these works function methodologically, much more so than aesthetically. although they are planned to carry a certain aesthetic viewpoint and weight. Ultimately, they are sites where Speculocultural Technopoiesis operates in practice. A space where unsettling technical assumptions, remixing computational elements, embodying cultural knowledge, and transmitting procedural insights occurs through sustained cycles of making and reflection. The modifications developed through these works are not one-off artistic decisions but transmissible procedures: approaches to spatialization that encode migration and memory, sonification strategies that honor material transformation as ritual, synthesis architectures that center marginalized bodies as sources of creative agency. These procedures, documented through code, video, reflective writing, and the works themselves, constitute the primary knowledge contribution of this research.

What these works collectively reveal is that modification is far, far from being simple technical problem-solving. Modification arises from a desire, a need, and is itself a form of cultural resistance and epistemological intervention. When I reconfigure a wavetable synthesizer to respond to hair texture, when I develop spatialization algorithms informed by forced migration patterns, when I mix hydrophone recordings with cosmic sonifications to connect historical trauma and speculative futures, I am insisting that different cultural logics produce different technological possibilities, and that the work of making those possibilities real is itself a form of scholarship. These sonic sculptures are arguments made through practice, demonstrations that research-creation means allowing the work itself to generate the questions, methods, and insights that conventional academic frameworks all too often predetermine.